\section{What Can the Selector Observe?}
\label{sec:diagnostics}

The theory identifies predictor response as the carrier of candidate-specific
utility. A selector still needs to know whether that signal is strong enough to
support a decision. We use three checks that move from the routing model itself
to the geometry of the candidate pool.

\subsection{Screening signal}

The cached stage should preserve candidates that matter. We compare its score
with the true marginal utility of the frozen predictor and report both rank
correlation and shortlist recall. A random ranking provides a simple reference:
a shortlist of size $q$ retains the best of $K$ candidates with probability
$q/K$. When the observed recall stays near this level, the response-aware stage
is being asked to rerank a shortlist that already lost much of the available
value.

\subsection{Within-state rankability}

Selection depends on the ordering of candidates inside one decision state. We
therefore fit diagnostic models to the true marginals using the information
available before the target is known: anchor features, candidate features,
candidate identity, the predictor response and simple response interactions.
We evaluate held-out $R^2$ and within-state rank correlation. An in-sample fit
and a label-permutation control check that the diagnostic pipeline itself can
learn when a signal is present. This test asks a direct question: can the
candidate ordering be recovered from the observations available to a router?

\subsection{Candidate novelty}

A cheaper signal comes from the candidate pool itself. For each candidate
series, we measure the share of its variance that remains after projecting it on
the anchor,
\begin{equation}
 \nu_j =
 \frac{\mathrm{Var}\!\left(c_j-
 \mathrm{proj}_{\mathrm{anchor}}(c_j)\right)}
 {\mathrm{Var}(c_j)} .
 \label{eq:novelty}
\end{equation}
Low novelty means that the candidate mostly repeats information already present
in the anchor. High novelty means that it contributes a larger independent
component. The statistic uses the training split and no target labels at
decision time. We fit any threshold on held-out targets within the deployment;
we do not assume that one threshold transfers across datasets.

\subsection{From observation to action}

These checks serve different purposes. Screening correlation tells us whether a
two-stage router can preserve useful candidates. Within-state rankability asks
whether the available observations contain enough information to order them.
Novelty describes a simpler deployment choice: whether relevance-based
selection is likely to add information or mostly repeat the anchor.

Structural headroom remains useful as an opportunity measure, but it does not
by itself predict how much a learned router will realize. Across the six traffic
benchmarks, targets with similar oracle headroom can have very different gains
over pooling. The experiments show that the decisive quantity is the fraction
of candidate utility that is recoverable under the deployment's own predictor,
pool and decision rule.
